\chapter{Forms and input validation}
\label{ch:forms}

\begin{aside}
A form is a list of inputs whose values must satisfy some
constraints before the user can submit. The constraints become
visible feedback; the inputs need shared state.
\end{aside}

\section{The form model}

A form is a struct of signals:

\begin{lstlisting}
struct LoginForm {
    Signal<std::string> username;
    Signal<std::string> password;
    Signal<std::string> error;
};
\end{lstlisting}

Inputs bind to fields; validation runs on change.

\section{Validation rules}

A rule is a function from value to optional error message:

\begin{lstlisting}
auto required(std::string_view v)
    -> std::optional<std::string>
{
    if (v.empty()) return "required";
    return std::nullopt;
}

auto min_length(int n) {
    return [n](std::string_view v)
        -> std::optional<std::string>
    {
        return v.size() >= n
             ? std::nullopt
             : std::optional<std::string>("too short");
    };
}
\end{lstlisting}

\section{Composing rules}

Multiple rules per field:

\begin{lstlisting}
auto validate_password = [&](std::string_view v) {
    if (auto e = required(v)) return e;
    if (auto e = min_length(8)(v)) return e;
    return std::optional<std::string>{};
};
\end{lstlisting}

\section{Showing errors}

Bind error display to the field's validation result:

\begin{lstlisting}
auto password_field = column(
    text_input(form.password),
    text(validate_password(ctx.get(form.password)).value_or(""))
        | fg<theme::error>
);
\end{lstlisting}

The error appears beneath the field, in red.

\section{Submit gating}

The submit button is disabled until all fields pass:

\begin{lstlisting}
auto can_submit = memo([&] {
    return !validate_username(ctx.get(form.username))
        && !validate_password(ctx.get(form.password));
});

auto button = ctx.get(can_submit)
    ? submit_button()
    : submit_button() | disabled;
\end{lstlisting}

\section{Async validation}

Some validations need a server check (``username taken''):

\begin{lstlisting}
ctx.async_effect([&]() -> Task<void> {
    auto u = ctx.get(form.username);
    auto avail = co_await check_available(u);
    ctx.set(username_error,
        avail ? "" : "taken");
});
\end{lstlisting}

Debounce to avoid flooding the server.

\section{Form-level errors}

Errors that don't belong to a single field (login failed)
live in a top-level error signal, shown above the form.

\section{Reset and dirty tracking}

Track which fields the user has touched; show errors only
after touch (don't shout at an empty fresh form).

\begin{lstlisting}
struct Field {
    Signal<std::string> value;
    Signal<bool> touched;
};
\end{lstlisting}

\section{Recap}

\begin{recap}
\begin{itemize}[leftmargin=*]
  \item Form = struct of signals.
  \item Rules return optional error messages.
  \item Submit gated on validation passing.
  \item Async validation for server checks.
  \item Touch tracking to avoid premature errors.
\end{itemize}
\end{recap}

\section*{Workshop}

\begin{enumerate}[leftmargin=*]
  \item Build a login form with username and password.
  \item Add an async ``username available'' check.
  \item Implement touch tracking so errors appear only
    after blur.
\end{enumerate}
